\documentclass[conference]{IEEEtran}
\IEEEoverridecommandlockouts
% The preceding line is only needed to identify funding in the first footnote. If that is unneeded, please comment it out.
\usepackage{cite}
\usepackage{amsmath,amssymb,amsfonts}
\usepackage{algorithmic}
\usepackage{graphicx}
\usepackage{textcomp}
\usepackage{booktabs}
\usepackage{multirow}
\usepackage{xcolor}
\def\BibTeX{{\rm B\kern-.05em{\sc i\kern-.025em b}\kern-.08em
    T\kern-.1667em\lower.7ex\hbox{E}\kern-.125emX}}
\begin{document}

\title{LSH Bucket-Occupancy Retention for Bounded Streaming Memory}


\author{\IEEEauthorblockN{1\textsuperscript{st} Chirantan Jain}
\IEEEauthorblockA{\textit{Computer Science Department} \\
\textit{New York University}\\
New York, USA \\
cj2735@nyu.edu}
\and
\IEEEauthorblockN{2\textsuperscript{nd} Ananya Agarwal}
\IEEEauthorblockA{\textit{Computer Science Department} \\
\textit{New York University}\\
New York, USA \\
aa13549@nyu.edu}
}

\maketitle

\begin{abstract}
Language models, AI agents, and streaming analytics systems operate
under strictly bounded memory. When ingesting an unbounded stream of
text embeddings, a retention policy must decide what to keep and what to
evict so that the retained buffer maximizes semantic diversity and
retrieval recall. Existing approaches either ignore content (FIFO,
Reservoir Sampling) or require expensive per-item scans of the full
retained set (exact semantic deduplication). We propose Bucket-Occupancy
Retention, a method that uses Locality-Sensitive Hashing (LSH) bucket
counts as a constant-time redundancy signal. Each incoming embedding is
hashed into L independent tables and scored by the median occupancy of
its buckets; a max-heap evicts the most redundant item when the budget
is full. On a 1M-passage MS MARCO stream with 50 percent injected
duplicates, the method achieves perfect Recall at 10 at half the memory
budget while all baselines plateau below 0.70. On the LoCoMo
conversational QA benchmark, it improves LLM question-answering accuracy
by 52 percent over a sliding window at a 25 percent memory budget.
Pairing the policy with Johnson-Lindenstrauss projection further
multiplies recall by 4x under tight memory by trading vector fidelity
for item capacity, though a distortion ceiling limits gains beyond 15
percent budget. We also report an honest negative result: under pure
topic drift without duplication, FIFO outperforms the proposed method,
establishing clear applicability boundaries.
\end{abstract}

\begin{IEEEkeywords}
locality-sensitive hashing, streaming memory, retention policy, semantic
deduplication, bounded memory, embedding streams
\end{IEEEkeywords}

\section{Introduction}

Modern language models, AI agents, and streaming analytics systems share
a fundamental systems constraint: context windows and fast-memory
buffers are strictly bounded in size. When an unbounded stream of text
embeddings $\{x_1, x_2, \ldots\}$ arrives one item at a time, a
\emph{retention policy} must maintain a memory buffer $M_t$ subject to
$|M_t| \le B$ while maximizing the semantic diversity and retrieval
recall of the retained set for arbitrary future queries---without
knowing those queries in advance.

Three families of policies dominate practice, each with a critical
limitation. \textit{Recency-based} methods (e.g., FIFO sliding windows)
discard the oldest items regardless of semantic value, catastrophically
forgetting rare but important early context. \textit{Uniform sampling}
methods (e.g., Reservoir Sampling~\cite{b10}) provide unbiased coverage
but waste capacity on redundant copies---in a stream with 50\%
duplicates, roughly half the buffer is squandered. \textit{Exact
semantic deduplication}~\cite{b11} performs nearest-neighbor search
against the full retained set for every incoming item, achieving high
accuracy but at $O(|M|)$ per-insertion cost that is prohibitive at
scale: our measurements show p99 latencies of 2,400--7,000~$\mu$s
compared with 27--35~$\mu$s for our method.

The gap is clear: no existing policy is simultaneously
\emph{semantically aware}---actively suppressing redundancy to preserve
diversity---and \emph{sublinear} in retained-set size at insertion time.

We address this gap with \textbf{Bucket-Occupancy Retention}, a method
that repurposes Locality-Sensitive Hashing (LSH) bucket counts as a
constant-time novelty signal. Each incoming embedding is hashed through
$L$ independent random-hyperplane tables; its \emph{redundancy score}
is the median occupancy of the buckets it lands in. Items in sparse
(novel) regions score low and are kept; items in dense (redundant)
regions score high and are evicted via a max-heap. The entire decision
costs $O(LKd + \log B)$ per item, independent of the retained-set size.

We evaluate the method on three synthetic stream scenarios derived from
the MS MARCO passage corpus~\cite{b17} (1M real web passages, 384-d
embeddings), plus a downstream LLM question-answering benchmark
(LoCoMo). Our key findings are:

\begin{enumerate}
\item \textbf{Perfect recall under duplication.}
      On a 50\%-duplicate stream, BucketOccupancy achieves
      Recall@10~=~1.0 at 50\% memory budget, while Reservoir and FIFO
      plateau at $\sim$0.69 (Section~V-B).
\item \textbf{Strong gains on clean data.}
      Even without injected duplicates, the policy exploits natural
      cluster imbalance to achieve $10.9\times$ higher Recall@10 than
      FIFO at 1\% budget and near-perfect Recall@1 ($\ge$0.994) across
      all budgets (Section~V-B).
\item \textbf{Downstream reasoning improvement.}
      At 25\% memory budget, BucketOccupancy improves LLM-QA accuracy
      by 52.3\% over FIFO on the LoCoMo benchmark by preserving
      semantic milestones that a sliding window forgets
      (Section~V-D).
\item \textbf{JL projection compounds the advantage.}
      Pairing the policy with Johnson-Lindenstrauss projection from
      $d{=}384$ to $d{=}64$ stores $6\times$ more items under the same
      byte budget, boosting recall by $4\times$ at tight budgets---until
      a distortion ceiling caps further gains (Section~V-F).
\item \textbf{Honest negative result.}
      Under pure topic drift without duplication, FIFO outperforms
      BucketOccupancy at every budget, establishing a clear boundary
      condition for diversity-first retention (Section~V-B).
\end{enumerate}

The remainder of the paper is organized as follows.
Section~II reviews related work on streaming retention, LSH, and
high-dimensional geometry.
Section~III presents the theoretical methodology, including the formal
problem statement, algorithm design, complexity analysis, and
concentration-inequality framing.
Section~IV describes the implementation: dataset preparation, evaluation
framework, and experimental environment.
Section~V reports results across all scenarios and ablations.
Section~VI concludes with deployment guidance, limitations, and future
work.

\section{Literature Review}

\subsection{Problem Setting: Bounded Streaming Memory}
Retention under bounded memory is a classical streaming problem.
Muthukrishnan's survey~\cite{b8} formalizes the one-pass constraint: an
algorithm may store at most $B$ items from a potentially infinite stream and
must make irrevocable keep/discard decisions. Within this framework, three
broad families of policies have been studied.

\textit{Recency-based} methods such as sliding-window FIFO buffers retain
only the $B$ most recent items. These are optimal when query relevance
is temporally local and the data distribution is stationary within the
window~\cite{b9}. However, they are blind to semantic content and will
overwrite rare but important early items as the stream advances.

\textit{Uniform sampling} methods, notably Vitter's Algorithm~R
(Reservoir Sampling)~\cite{b10}, guarantee that every stream element has
equal probability $B/N$ of appearing in the final sample. This provides
unbiased global coverage but makes no attempt to suppress redundancy; in
a stream with 50\% exact duplicates, roughly half the retained set will
be wasted on duplicate copies.

\textit{Exact semantic deduplication} methods perform nearest-neighbor
lookup against the full retained set for every incoming item. Abbas et
al.~\cite{b11} demonstrate this approach for large-scale dataset pruning
using FAISS, discarding any item whose cosine similarity to an existing
member exceeds $1-\varepsilon$. While highly accurate, the per-insertion
cost scales as $O(|M|)$ in practice, making it prohibitive for
high-throughput streaming: our experiments show that p99 latency reaches
2,400--4,700~$\mu$s at moderate budgets, compared with 29--35~$\mu$s
for our proposed method.

The gap that motivates this work is the absence of a policy that is
simultaneously \emph{semantically aware} (actively suppresses redundancy)
and \emph{sublinear} in retained-set size at insertion time. Our
Bucket-Occupancy Retention method addresses this gap using LSH-based
density estimation.

\subsection{Baseline Taxonomy and Assumptions}
To position our contribution, we implemented four baselines that span the
families above:
\begin{itemize}
\item \textbf{FIFO (First-In, First-Out):} a sliding window of size $B$
      that assumes recency equals relevance.
\item \textbf{Reservoir Sampling}~\cite{b10}: Vitter's Algorithm~R,
      providing a uniform random sample of size $B$ over the full stream.
\item \textbf{Stream-LSH}~\cite{b12}: a scoring method that ranks
      retained items by a weighted combination of freshness (exponential
      time-decay), LSH collision penalties, and an external per-item
      quality signal (e.g., click-through rate). In our ingestion-only
      setting no external quality signal is available, so the quality
      term is zero; the LSH collision penalty is small relative to the
      exponential decay, causing the composite score to be dominated by
      recency. Empirically, Stream-LSH produces identical recall to
      FIFO at every budget we tested.
\item \textbf{Semantic Dedup (SOTA Gold Standard)}~\cite{b11}: exact
      nearest-neighbor cosine-similarity filtering via FAISS. Serves as
      an accuracy ceiling but is computationally infeasible for fast
      streaming (restricted to $B/N \le 0.05$ in our experiments due to
      latency).
\end{itemize}
Our end-to-end evaluation across three synthetic stream scenarios
(CleanStream, HeavyDuplication, TopicDrift) reveals that the optimal
baseline varies by regime, confirming that no single assumption
(recency, randomness, or exact matching) dominates universally.

\subsection{Locality-Sensitive Hashing and Novelty Estimation}
Locality-Sensitive Hashing was introduced by Indyk and
Motwani~\cite{b13} to solve approximate nearest-neighbor search in
sublinear time. The foundational property is that for a hash family
$\mathcal{H}$ and similarity measure $\mathrm{sim}(\cdot,\cdot)$:
\begin{equation}
\Pr_{h \sim \mathcal{H}}[h(x)=h(y)] = f\!\bigl(\mathrm{sim}(x,y)\bigr),
\label{eq:lsh-collision}
\end{equation}
where $f$ is monotonically increasing. For random-hyperplane (SimHash)
projections applied to unit-normalized vectors, Charikar~\cite{b14}
showed that $\Pr[h(x)=h(y)] = 1 - \theta(x,y)/\pi$, where $\theta$ is
the angle between $x$ and $y$.

Our method repurposes this collision-probability relationship: rather
than using LSH to retrieve neighbors, we use bucket \emph{occupancy
counts} across $L$ independent hash tables as a streaming proxy for
local semantic density. High occupancy signals a crowded (redundant)
region; low occupancy signals a sparse (novel) one. This reframing
avoids any scan over the retained set, reducing novelty estimation to
$O(L)$ hash lookups.

We additionally explore a Jaccard-over-signatures variant, where the
redundancy score is defined as the maximum Jaccard similarity of the
incoming item's $L$-element bucket signature against retained items
sharing at least one bucket. An inverted index keyed by
$(\text{table},\text{bucket})$ pairs limits candidate retrieval to
$O(L \times \bar{b})$, where $\bar{b}$ is the average bucket
population, rather than scanning the full retained set.

\subsection{High-Dimensional Geometry and Distance Concentration}
In high-dimensional spaces, random vectors exhibit distance
concentration: pairwise distances cluster tightly around the
mean~\cite{b15}. This phenomenon explains our empirical bucket
histograms: retained-set occupancy distributions are sharply peaked
rather than long-tailed, because most items are roughly equidistant from
one another. Blum et al.~\cite{b15} provide the formal treatment showing
that for random unit vectors in $\mathbb{R}^d$, the standard deviation
of pairwise distance shrinks as $O(1/\sqrt{d})$.

This has two practical implications for our work. First, it explains why
BucketOccupancy performs comparably to (rather than dramatically better
than) baselines on TopicDrift, where the stream lacks strong redundancy
structure. Second, it motivates the Johnson-Lindenstrauss
experiments (Section~V-F): JL projection from $d{=}384$ to $d{=}64$
preserves pairwise distances within a $(1{\pm}\varepsilon)$ factor with
high probability~\cite{b16}, enabling 6$\times$ higher item capacity
under the same memory footprint.

\subsection{Known Limitations and Honest Positioning}
A persistent finding in the streaming selection literature is
\emph{objective mismatch}: maximizing global semantic diversity is not
always equivalent to maximizing query-time relevance, particularly
under temporal non-stationarity~\cite{b9}. Our TopicDrift experiments
confirm this: FIFO outperforms BucketOccupancy at every budget level
(e.g., R@10 of 0.2696 vs.\ 0.2333 at 25\% budget) when the query
distribution is biased toward recent topics.

We present this as a principled \emph{boundary condition} of
diversity-first retention, not an anomaly. In streams where redundancy
and duplication are the dominant dynamics---which is the common case in
chat histories, sensor logs, and repeated prompt contexts---our method
provides its strongest advantage.

%------------------------------------------------------------------------
\section{Theoretical Methodology}

\subsection{Formal Problem Statement}
Let $\{x_t\}_{t \ge 1}$ be an unbounded stream of $d$-dimensional
embedding vectors arriving one at a time. At each time step $t$, a
retention policy maintains a memory buffer $M_t$ subject to a hard
budget constraint $|M_t| \le B$. The policy must make an irrevocable
binary decision---keep or discard---for every incoming $x_t$.

The optimization objective is to maximize the \emph{semantic coverage}
of $M_t$: informally, the retained set should contain at least one
representative close to every distinct region of the embedding space
visited by the stream. We operationalize this via Recall@$k$, cosine diversity metrics,
and downstream QA accuracy (Section~V).

\subsection{Algorithm Design and Decision Rule}
\subsubsection{LSH Signature Computation}
Each incoming embedding $x_t \in \mathbb{R}^d$ is hashed through $L$
independent random-hyperplane hash functions, each producing a $K$-bit
signature. Concretely, let $\{r_{\ell,1},\ldots,r_{\ell,K}\}$ be
i.i.d.\ $\mathcal{N}(0,I_d)$ projection vectors for table $\ell$. The
bucket ID for table $\ell$ is:
\begin{equation}
h_\ell(x_t) = \sum_{j=1}^{K} \mathbf{1}[\langle r_{\ell,j}, x_t \rangle > 0] \cdot 2^{K-j}.
\label{eq:hash}
\end{equation}
This yields an $L$-element \emph{bucket signature}
$\sigma(x_t) = (h_1(x_t),\ldots,h_L(x_t))$.

\subsubsection{Redundancy Scoring}
We maintain per-table bucket counts $c_\ell(b)$ tracking how many
currently retained items hash to bucket $b$ in table $\ell$. All counts
are initialized to zero and updated incrementally: $c_\ell(b)$ is
incremented when an item hashing to bucket $b$ in table $\ell$ is
admitted, and decremented when such an item is evicted. Two
aggregation strategies are supported:

\textit{Median-occupancy} (default):
\begin{equation}
s(x_t) = \mathrm{median}\bigl(c_1(h_1(x_t)),\; \ldots,\; c_L(h_L(x_t))\bigr).
\label{eq:median-occ}
\end{equation}

\textit{Max-Jaccard}: for each candidate $y \in M_t$ sharing at least
one bucket with $x_t$ (retrieved via an inverted index), compute the
signature-level Jaccard similarity. Each item's signature is the set of
$L$ bucket IDs it maps to, one per table; since both sets have
cardinality $L$, the union has size $2L - |\text{intersection}|$,
yielding:
\begin{equation}
J(x_t,y) = \frac{|\{\ell : h_\ell(x_t) = h_\ell(y)\}|}{2L - |\{\ell : h_\ell(x_t) = h_\ell(y)\}|},
\label{eq:jaccard}
\end{equation}
and set $s(x_t) = \max_{y \in \mathcal{C}(x_t)} J(x_t,y)$, where
$\mathcal{C}(x_t)$ is the candidate set from the inverted index.

In both variants, a \emph{lower} score indicates greater novelty.

\subsubsection{Capacity-Bounded Eviction}
When $|M_t| < B$, every item is accepted. Once the buffer is full, a
new item $x_t$ is admitted only if $s(x_t) < s(y^*)$, where $y^*$ is
the most redundant item in $M_t$ (maintained at the top of a max-heap
keyed on redundancy score). If admitted, $y^*$ is evicted and all
bucket counts are updated accordingly. This one-in-one-out mechanism
is detailed in Algorithm~1.

\begin{figure}[h]
\centering
\fbox{\parbox{0.92\columnwidth}{
\textbf{Algorithm 1:} BucketOccupancy Insert$(x_t)$\\[4pt]
1: Compute signature $\sigma(x_t) = (h_1(x_t),\ldots,h_L(x_t))$\\
2: Compute redundancy score $s(x_t)$ via~(\ref{eq:median-occ}) or~(\ref{eq:jaccard})\\
3: \textbf{if} $|M_t| < B$ \textbf{then}\\
4: \quad Add $x_t$ to $M_t$; push $(x_t, s(x_t))$ to max-heap\\
5: \quad $\forall\,\ell$: $c_\ell(h_\ell(x_t)) \mathrel{+}= 1$ \hfill\textit{// increment bucket counts}\\
6: \textbf{else}\\
7: \quad Let $y^* = \arg\max_{y \in M_t} s(y)$ \hfill (heap top)\\
8: \quad \textbf{if} $s(x_t) < s(y^*)$ \textbf{then}\\
9: \quad\quad Evict $y^*$; $\forall\,\ell$: $c_\ell(h_\ell(y^*)) \mathrel{-}= 1$\\
10: \quad\quad Add $x_t$ to $M_t$; push $(x_t, s(x_t))$ to max-heap\\
11: \quad\quad $\forall\,\ell$: $c_\ell(h_\ell(x_t)) \mathrel{+}= 1$\\
12: \quad \textbf{else} Discard $x_t$\\
13: \textbf{end if}
}}
\label{alg:insert}
\end{figure}

\subsection{Complexity Analysis}
\label{sec:complexity}

\subsubsection{Per-Insertion Cost}
The dominant operations per arriving item are:
\begin{itemize}
\item \textbf{Hashing:} $L \times K$ inner products of dimension $d$,
      yielding $O(LKd)$. With our defaults ($L{=}8$, $K{=}10$,
      $d{=}384$), this is $\approx$30,720 multiply-adds.
\item \textbf{Occupancy aggregation:} $O(L)$ hash-table lookups for
      median mode. For Jaccard mode, $O(L \cdot \bar{b})$ where
      $\bar{b} = |M_t|/2^K$ is the average number of items per bucket.
\item \textbf{Heap maintenance:} $O(\log B)$ for push/pop when the
      buffer is full.
\end{itemize}
Total per-item cost is therefore $O(LKd + \log B)$ for median mode,
independent of $|M_t|$. This is the critical theoretical advantage over
exact semantic deduplication, which requires $O(|M_t| \cdot d)$ per
insertion for a brute-force scan (or $O(d \log |M_t|)$ with
tree-based indices that are difficult to maintain under eviction).

\subsubsection{Space Overhead}
Beyond the $B$ retained embeddings (which every policy must store),
BucketOccupancy adds three small auxiliary structures: the random
projection vectors used for hashing ($L \times K \times d$ floats
$\approx$ 120~KB), a fixed-size table of bucket counters
($L \times 2^K$ integers $\approx$ 32~KB), and a max-heap of $B$
entries for eviction. In total, these add $<$\textbf{1\%} to the
memory footprint of the embeddings themselves (e.g., $\sim$1~MB of
overhead on 153.6~MB of embeddings at $B{=}100{,}000$). By contrast,
Semantic Dedup requires a FAISS nearest-neighbor index that grows
linearly with $|M|$.

\subsubsection{Empirical Throughput Validation}
Our throughput measurements confirm the complexity analysis:
BucketOccupancy sustains $>$75,000 items/sec with median latency of
29--35~$\mu$s, compared to Semantic Dedup's p99 latency of
2,400--4,700~$\mu$s at budgets as low as 5\%. Reservoir Sampling, which
performs no semantic computation, runs at $\sim$950,000 items/sec---roughly
12$\times$ faster---but without any redundancy awareness.

\subsection{Collision Probability and the Indyk-Motwani Connection}
\label{sec:indyk-motwani}
The theoretical justification for using bucket occupancy as a
redundancy signal derives from the LSH collision-probability
guarantee~(\ref{eq:lsh-collision}). For random-hyperplane
hashing~\cite{b14}, items that are semantically similar (high cosine
similarity) collide with high probability across multiple tables. This
means:
\begin{itemize}
\item An item landing in \emph{dense} buckets has many near-duplicates
      already retained $\Rightarrow$ high $s(x_t)$ $\Rightarrow$
      redundant.
\item An item landing in \emph{sparse} buckets occupies a region not
      well-represented in $M_t$ $\Rightarrow$ low $s(x_t)$
      $\Rightarrow$ novel.
\end{itemize}

Our internal characterization experiments (Section~V-A) provide
empirical consistency checks for this theoretical link:
\begin{enumerate}
\item \textbf{Duplicate filtering:} When exact duplicates are injected
      at rates of 10\%--50\%, the policy identifies and discards them at
      $>$90\% accuracy, confirming that high-similarity items reliably
      produce high occupancy scores.
\item \textbf{Occupancy flattening:} Under capacity-bounded eviction,
      the retained-set bucket histogram shifts from a long-tailed
      distribution (reflecting natural cluster imbalance) to a tightly
      peaked one, demonstrating active suppression of the densest
      semantic regions.
\item \textbf{Hash-granularity stability:} Sweeping $K$ from 6 to 14
      yields stable Recall@10 performance, indicating that the
      occupancy signal is robust to moderate changes in hash resolution.
      $K{=}10$ provides a practical balance for 384-dimensional
      embeddings.
\end{enumerate}

\subsection{Concentration Inequality Framing}
\label{sec:concentration}
The median-occupancy score $s(x_t)$ aggregates $L$ independent bucket
counts. Since each table uses an independent random projection, the
per-table counts $c_\ell(h_\ell(x_t))$ are independent random variables
(conditioned on $x_t$ and the current retained set). By Chebyshev's
inequality:
\begin{equation}
\Pr\!\left[\,|s(x_t) - \mathbb{E}[s(x_t)]| \ge \delta\,\right]
\;\le\; \frac{\mathrm{Var}[s(x_t)]}{\delta^2},
\label{eq:chebyshev}
\end{equation}
where the variance decreases as $L$ increases, since the median of $L$
independent samples concentrates around the population median.

In practice, this means the retention decision is \emph{stable} across
random seeds and hash-function draws: an item that is truly redundant
will consistently receive a high score, and a truly novel item will
consistently receive a low score. Our $K$-sweep results (stable recall
across $K{=}6$--$14$ with $L{=}8$) provide indirect empirical support for
this concentration behavior.

\subsection{Theory--Result Alignment Across Experimental Scenarios}
\label{sec:alignment}

Each major experimental result can be connected to a theoretical
expectation:

\subsubsection{CleanStream (No Duplicates)}
Even without injected duplicates, real embedding distributions exhibit
natural cluster imbalance. The occupancy-based policy exploits this
structure by ensuring every semantic region has at least one
representative. At 1\% budget, BucketOccupancy achieves R@10~=~0.1083
(vs.\ 0.0099 for Reservoir, a 10$\times$ improvement) and near-perfect
R@1~$\ge$~0.994 across all budgets, confirming that the policy
prioritizes retaining the single nearest neighbor to each query region.

\subsubsection{Heavy Duplication (50\% Duplicates)}
This is the regime where the redundancy signal is strongest, and the
theoretical prediction is correspondingly strong. At 50\% budget,
BucketOccupancy achieves \emph{perfect recall} (R@10~=~1.0) while
Reservoir and FIFO plateau at $\sim$0.69. The policy ``sees through''
the duplication: by evicting redundant copies via the max-heap, it
retains all 1M unique items within just half the total stream budget.
At 25\% budget, BucketOccupancy leads Reservoir by +9 percentage points
(0.4992 vs.\ 0.4080).

\subsubsection{Topic Drift}
Under sequential topic drift without heavy duplication, the
diversity-first objective mismatches the recency-biased query
distribution. FIFO leads at every budget (e.g., R@10~=~0.2696 vs.\
0.2333 at 25\%). This is theoretically expected: when relevance is
temporally concentrated, a recency prior outperforms a global diversity
prior. Notably, BucketOccupancy still produces a more diverse retained
set (intra-set distance 1.025 vs.\ 0.901 for FIFO at 25\%), but that
diversity does not translate to recall under temporal query bias.

\subsubsection{LoCoMo LLM-QA Benchmark}
Downstream question-answering on the LoCoMo benchmark (10 multi-session
conversations, Qwen~2.5~7B) confirms that retained semantic diversity
translates to reasoning performance. At 25\% budget, BucketOccupancy
achieves 7.37\% QA accuracy vs.\ 4.84\% for FIFO (+52\% relative
gain). Even at 50\% budget, BucketOccupancy maintains 72\% of the
full-memory accuracy (11.60\% / 16.04\%), demonstrating graceful
degradation under memory pressure.

\subsubsection{Johnson-Lindenstrauss Projection Tradeoff}
JL projection from $d{=}384$ to $d{=}64$ enables storing 6$\times$ more
items under the same byte budget. At 5\% memory footprint,
BucketOccupancy+JL achieves R@10~=~0.3892 vs.\ 0.0989 for
uncompressed BucketOccupancy (a 4$\times$ improvement). However,
compression distortion imposes a hard ceiling: $d{=}64$ recall flatlines at
0.5561 regardless of further capacity increases. At 50\% budget,
uncompressed 384D vectors achieve near-perfect recall (0.9999),
dominating the compressed variants. This Pareto tradeoff---lossy
compression wins at tight budgets, full fidelity wins at loose
budgets---is precisely the behavior predicted by JL theory~\cite{b16}.

\subsection{Threats to Validity}
\label{sec:threats}
We identify the following limitations of our theoretical and empirical
framework:
\begin{itemize}
\item \textbf{Embedding quality:} All results assume a fixed,
      pre-trained embedding model (384-dimensional). The occupancy
      signal's effectiveness depends on the embedding space faithfully
      capturing semantic similarity.
\item \textbf{Stationarity assumption:} The collision-probability
      argument assumes the hash functions are drawn independently of the
      data. Adversarial or highly non-stationary streams may violate
      this.
\item \textbf{Hyperparameter sensitivity:} While the $K$-sweep shows
      robustness over $K \in [6,14]$, the interaction between $L$, $K$,
      and the intrinsic dimensionality of the data is not fully
      characterized.
\item \textbf{Oracle construction:} Recall metrics depend on the oracle
      query set and its alignment with the stream. Different query
      distributions could change the relative ranking of policies.
\item \textbf{Jaccard overhead:} The inverted-index variant introduces
      candidate-dependent cost $O(L \cdot \bar{b})$ that can degrade
      under skewed distributions with very dense buckets.
\end{itemize}
These caveats motivate future work on adaptive hash functions (e.g.,
learned projections, online gradient methods) and hybrid
recency-diversity policies that can handle both redundancy and drift.

%------------------------------------------------------------------------
\section{Implementation}

\subsection{Dataset Preparation}
\textbf{Source data.}
We use MS MARCO v2.1~\cite{b17}, a corpus of 1,000,000 real web
passages from Microsoft's Bing search logs.

\textbf{Embedding model.}
All passages are embedded with \texttt{all-MiniLM-L6-v2}
(a 6-layer distilled Sentence-BERT model from the
Sentence-Transformers library), producing
384-dimensional L2-normalized vectors.

\textbf{Stream variants.}
Three synthetic streams stress-test different retention regimes:

\begin{table}[h]
\centering
\caption{Stream variants used in evaluation.}
\label{tab:streams}
\begin{tabular}{@{}llp{4.2cm}@{}}
\toprule
Stream & Size & Description \\
\midrule
\texttt{CleanStream\_0}        & 1M & Original passages, no duplicates \\
\texttt{HeavyDup\_50}          & 2M & 1M originals + 1M exact copies (randomly selected sources), shuffled together \\
\texttt{TopicDrift}            & 1M & Sorted by 20 clusters (MiniBatch $K$-Means on the 384-d embeddings) to simulate sequential topic shift \\
\bottomrule
\end{tabular}
\end{table}

\textbf{Queries and oracle.}
5,000 passages are randomly sampled \emph{from the same 1M-passage
corpus} to serve as queries; because each query is itself a stream item,
it has a trivially perfect nearest neighbor (itself) plus 99 additional
close neighbors.
Oracle ground truth is the exact top-100 nearest neighbors for each
query computed via brute-force inner product
(FAISS \texttt{IndexFlatIP}) over the full unfiltered stream.
Recall@$k$ then measures what fraction of each query's oracle top-$k$
neighbors the retention policy managed to keep in its buffer.

\subsection{Evaluation Framework}
The pipeline operates in four stages: (1)~stream all $N$ items through
the policy one-by-one, recording per-item nanosecond latency;
(2)~extract the retained set $(M, \text{kept\_ids})$;
(3)~build a FAISS \texttt{IndexFlatIP} index on $M$ and search with
held-out queries; (4)~compare retrieved items against the oracle ground
truth.

\begin{table}[h]
\centering
\caption{Evaluation metrics organized by tier. Tier~1 metrics are
reported in all experiments; Tier~2 metrics appear in selected
analyses where diversity structure is the focus (Sections~V-B, V-C);
Tier~3 is reported in the JL projection analysis (Section~V-F).}
\label{tab:metrics}
\begin{tabular}{@{}clp{3.8cm}@{}}
\toprule
Tier & Metric & Measures \\
\midrule
1 & Recall@$k$ ($k{=}1,10,100$) & Fraction of oracle top-$k$ items retained \\
1 & Latency (p50, p99, p999) & Per-insertion time distribution \\
2 & Mean Intra-Set Distance & Average pairwise spread of retained items \\
2 & Cosine Diversity Score & $1 - \text{mean NN cosine sim.}$ within retained set \\
2 & Cosine Coverage Radius & Worst-case angular gap to query set \\
3 & Memory Footprint (bytes) & Embeddings + hash tables + bucket counts \\
\bottomrule
\end{tabular}
\end{table}

\subsection{LoCoMo LLM-QA Evaluation}
To test downstream reasoning impact, we use the LoCoMo
benchmark---10 multi-session conversations with associated QA pairs.
Each conversation's turns are embedded, streamed through the retention
policy under budget constraints, and the top-5 most relevant retained
turns (by cosine similarity to the question embedding) are fed as
context to Qwen~2.5~7B, served locally via the Ollama inference
framework, to generate answers.
Scoring uses exact substring match as the primary criterion: a
prediction is correct if the ground-truth string appears as a substring
of the model's response, or vice versa. When neither direction matches,
the same Qwen~2.5~7B model is used as an LLM judge---prompted with
the ground-truth and predicted answers and asked to output ``YES'' or
``NO'' for semantic equivalence---providing a fuzzy-match fallback.

\subsection{Experimental Environment}
All streaming experiments run single-threaded for accurate latency
measurement. FAISS ground-truth computation uses an NVIDIA RTX~5070.
The platform is WSL2 (Ubuntu) on Windows with Python~3.x, FAISS,
FALCONN, scikit-learn, NumPy, Pandas, Seaborn, and Ollama.

%------------------------------------------------------------------------
\section{Results and Analysis}

\subsection{Internal Characterization}
Before comparing against baselines, we validated the algorithm's
internal mechanics through the following experiments, each targeting a
different aspect of the retention policy.

\subsubsection{Steady-State Memory Bounding}
With capacity set to 100,000, the retained-set size strictly flatlines
at the cap, seamlessly transitioning from ``accept everything'' to
one-in-one-out eviction.

\subsubsection{Bucket Distribution Flattening}
The original stream exhibits a long-tail bucket distribution. Under
capacity-bounded eviction, the policy actively equalizes bucket
occupancy, producing a tightly peaked distribution that confirms
uniform coverage of the embedding space.

\subsubsection{Duplicate Filtering Accuracy}

\begin{table}[h]
\centering
\caption{Duplicate filtering accuracy at varying injection rates.}
\label{tab:dup-filter}
\begin{tabular}{@{}rrrr@{}}
\toprule
Dup.\ Rate & Stream Size & Dups Discarded & Discard Rate \\
\midrule
10\% & 1,111,111 &  94,969 & 85.47\% \\
30\% & 1,428,571 & 379,085 & 88.45\% \\
50\% & 2,000,000 & 915,522 & 91.55\% \\
\bottomrule
\end{tabular}
\end{table}

The policy correctly identifies and discards injected duplicates at
rates exceeding 85\%, improving as duplication increases.

\subsubsection{Hash Granularity ($K$) Sweep}

\begin{table}[h]
\centering
\caption{Recall and throughput across hash lengths $K$ ($L{=}8$, $d{=}384$).
Each table has $2^K$ buckets.}
\label{tab:k-sweep}
\begin{tabular}{@{}rrrrr@{}}
\toprule
$K$ & Buckets/Table ($2^K$) & Recall@10 & Throughput (items/s) \\
\midrule
 6 &     64 & 0.1885 & 78,365 \\
 8 &    256 & 0.1895 & 77,251 \\
10 &  1,024 & 0.1886 & 74,604 \\
12 &  4,096 & 0.1878 & 69,007 \\
14 & 16,384 & 0.1865 & 60,153 \\
\bottomrule
\end{tabular}
\end{table}

Recall is stable across all $K$ values. $K{=}10$ provides the best
balance between hash quality and throughput for $d{=}384$.

\subsubsection{Distribution Drift Response}
Under sequential topic drift, the retention rate decays monotonically as
capacity is reached. The policy performs comparably to Reservoir
sampling, confirming that the bucket-novelty signal is most valuable
when redundancy---not drift---is the dominant structure.

%--------------------------------------------------------------------
\subsection{End-to-End Recall vs.\ Memory Budget}

\subsubsection{Clean Stream (No Duplicates)}

\begin{table}[h]
\centering
\caption{Recall on \texttt{CleanStream\_0} (1M items, no duplicates).}
\label{tab:clean}
\begin{tabular}{@{}rcccccc@{}}
\toprule
$B/N$ & \multicolumn{2}{c}{BucketOcc.} & Reservoir & FIFO & Stream-LSH \\
\cmidrule(lr){2-3}
      & R@10 & R@1 & R@10 & R@10 & R@10 \\
\midrule
 1\% & \textbf{0.1083} & \textbf{0.9940} & 0.0099 & 0.0088 & 0.0088 \\
 5\% & \textbf{0.1445} & \textbf{0.9988} & 0.0485 & 0.0448 & 0.0448 \\
10\% & \textbf{0.1886} & \textbf{0.9988} & 0.1003 & 0.0907 & 0.0907 \\
25\% & \textbf{0.3226} & \textbf{0.9990} & 0.2491 & 0.2257 & 0.2257 \\
50\% & \textbf{0.5447} & \textbf{0.9994} & 0.4984 & 0.4518 & 0.4518 \\
\bottomrule
\end{tabular}
\end{table}

BucketOccupancy significantly outperforms all baselines even on clean
data. At the 1\% budget it achieves $10.9\times$ higher R@10 than FIFO
(0.1083 vs.\ 0.0088). Its R@1 is near-perfect across all budgets
($\ge$0.994): because each query was sampled from the stream, its
top-1 oracle neighbor is itself, and BucketOccupancy almost always
retains it. The R@10 gap, which measures how many of the 10 nearest
neighbors survive, is where the diversity advantage is most visible.
This reveals that even ``clean'' data has natural semantic clustering,
and the policy exploits this structure by ensuring every region has at
least one representative.

Stream-LSH produces identical results to FIFO across all budgets,
confirming that without external quality signals its recency term
dominates. Semantic Deduplication achieves comparable recall
(R@10~=~0.1084 at 1\%) but was restricted to $B/N \le 0.05$ due to
its p99 latency of 2,435--7,051~$\mu$s vs.\ 27--31~$\mu$s for
BucketOccupancy. Because of this prohibitive cost, Semantic Dedup and
Stream-LSH are omitted from subsequent tables; they serve as an
accuracy ceiling and a FIFO-equivalent reference, respectively.

\textbf{Throughput:} Reservoir runs at $\sim$946,618 items/s vs.\
BucketOccupancy at $\sim$82,250 items/s ($\sim$11.5$\times$ slower).

\subsubsection{Heavy Duplication (50\% Duplicates)}

\begin{table}[h]
\centering
\caption{Recall on \texttt{HeavyDuplication\_50} ($N{=}2$M total items,
1M unique). Budget ratio $B/N$ is relative to the full stream size $N$;
e.g., $B/N{=}50\%$ means $B{=}1$M, which equals the number of unique items.}
\label{tab:heavy}
\begin{tabular}{@{}rccccc@{}}
\toprule
$B/N$ & \multicolumn{2}{c}{BucketOcc.} & Reservoir & FIFO \\
\cmidrule(lr){2-3}
      & R@10 & R@1 & R@10 & R@10 \\
\midrule
 1\% & 0.0193 & 0.0168 & 0.0190 & 0.0181 \\
 5\% & \textbf{0.0989} & 0.0830 & 0.0920 & 0.0912 \\
10\% & \textbf{0.1972} & 0.1790 & 0.1785 & 0.1783 \\
25\% & \textbf{0.4992} & \textbf{0.5144} & 0.4080 & 0.4035 \\
50\% & \textbf{1.0000} & \textbf{0.9994} & 0.6871 & 0.6832 \\
\bottomrule
\end{tabular}
\end{table}

This is the headline result. At 50\% budget on a 50\%-duplicate
stream, BucketOccupancy achieves \emph{perfect recall} (1.0000) while
all baselines plateau at $\sim$0.69. The policy ``sees through'' the
duplication: by evicting redundant copies, it retains all 1,000,000
unique items within exactly half the 2,000,000-item memory budget. At
25\% budget, BucketOccupancy leads Reservoir by +9.1 percentage points
(0.4992 vs.\ 0.4080).

\subsubsection{Topic Drift (Negative Result)}

\begin{table}[h]
\centering
\caption{Recall on \texttt{TopicDrift} (1M items, 20 sequential clusters).}
\label{tab:drift}
\begin{tabular}{@{}rcccc@{}}
\toprule
$B/N$ & BucketOcc.\ R@10 & Reservoir R@10 & FIFO R@10 \\
\midrule
 1\% & 0.0089 & 0.0099 & \textbf{0.0111} \\
 5\% & 0.0450 & 0.0525 & \textbf{0.0541} \\
10\% & 0.0888 & 0.1015 & \textbf{0.1030} \\
25\% & 0.2333 & 0.2479 & \textbf{0.2696} \\
50\% & 0.4817 & 0.5020 & \textbf{0.5224} \\
\bottomrule
\end{tabular}
\end{table}

Under topic drift without heavy duplication, BucketOccupancy performs
\emph{worse} than both FIFO and Reservoir at all budgets. FIFO leads
because recency correlates with relevance when topics shift
sequentially. This is an important \textbf{negative result}:
bucket-occupancy retention provides no advantage when distribution
drift is the dominant dynamic.

However, BucketOccupancy maintains a consistently higher mean intra-set
distance (1.025 vs.\ 0.901 for FIFO at 25\%), confirming it \emph{is}
producing a more diverse retained set---but that diversity does not
translate to recall when the query distribution is biased toward recent
topics.

%--------------------------------------------------------------------
\subsection{Cosine-Based Diversity Analysis}

We define two angular metrics over the retained set $M$:
\begin{itemize}
\item \textbf{Cosine Diversity Score:}
      $1 - \frac{1}{|M|}\sum_{x \in M}\max_{y \in M \setminus \{x\}}
      \cos(x,y)$. Higher values indicate that retained items are, on
      average, farther from their nearest retained neighbor---i.e., less
      internally redundant.
\item \textbf{Cosine Coverage Radius:}
      $\max_{q \in Q}\;\min_{x \in M}\;(1-\cos(q,x))$, where $Q$ is the
      held-out query set. This is the worst-case angular gap between any
      query and its closest representative in $M$; lower is better.
\end{itemize}

\begin{table}[h]
\centering
\caption{Cosine diversity metrics on \texttt{HeavyDuplication\_50}.}
\label{tab:cosine}
\begin{tabular}{@{}rlcc@{}}
\toprule
$B/N$ & Policy & Diversity Score & Coverage Radius \\
\midrule
 1\% & FIFO           & 0.5227 & 0.7678 \\
 1\% & Reservoir      & 0.5210 & 0.7746 \\
 1\% & \textbf{BucketOcc.} & \textbf{0.5299} & \textbf{0.7613} \\
\midrule
 5\% & FIFO           & 0.5215 & 0.6911 \\
 5\% & Reservoir      & 0.5253 & 0.7111 \\
 5\% & \textbf{BucketOcc.} & \textbf{0.5331} & \textbf{0.6876} \\
\midrule
50\% & FIFO           & 0.5246 & 0.6134 \\
50\% & Reservoir      & 0.5200 & 0.6151 \\
50\% & \textbf{BucketOcc.} & \textbf{0.5264} & \textbf{0.0000} \\
\bottomrule
\end{tabular}
\end{table}

BucketOccupancy achieves the highest Cosine Diversity Score across all
budgets and the lowest Cosine Coverage Radius (tighter angular
coverage). At the 50\% budget on the HeavyDuplication stream,
BucketOccupancy's coverage radius drops to effectively \textbf{zero}.
Because $B/N{=}50\%$ corresponds to $B{=}1$M items and the stream
contains exactly 1M unique vectors, the policy retains every unique
passage---so every query's nearest neighbor is present and the
worst-case angular gap vanishes. FIFO and Reservoir, which waste
capacity on duplicate copies, still leave coverage gaps of $\sim$0.61.

%--------------------------------------------------------------------
\subsection{Downstream Reasoning: LoCoMo LLM-QA}

\begin{table}[h]
\centering
\caption{Mean QA accuracy (\%) across 10 LoCoMo conversations.}
\label{tab:locomo}
\begin{tabular}{@{}rccc@{}}
\toprule
Budget & BucketOcc. & FIFO & Reservoir \\
\midrule
 10\% &  3.56 &  3.08 & \textbf{3.97} \\
 25\% & \textbf{7.37} &  4.84 &  6.49 \\
 50\% & \textbf{11.60} &  9.61 &  9.91 \\
100\% & \textbf{16.04} & 15.46 & 15.33 \\
\bottomrule
\end{tabular}
\end{table}

At the critical 25\% budget level, BucketOccupancy outperforms FIFO by
52.3\% relatively (7.37\% vs.\ 4.84\%). FIFO catastrophically forgets
early-conversation details needed for long-horizon questions. Even as
memory is halved from 100\% to 50\%, BucketOccupancy retains 72.3\% of
full-memory accuracy (11.60\%~/~16.04\%), demonstrating stable reasoning
under memory pressure.

At 10\% budget, Reservoir slightly outperforms both methods (3.97\%),
likely because at extreme compression the random diversity of Reservoir
provides marginally better coverage than BucketOccupancy's hash-based
decisions on very short conversation windows.

\textit{Note on absolute accuracy:} LoCoMo is a challenging benchmark
with multi-hop, temporal-reasoning questions over long conversations;
the absolute accuracy of a 7B-parameter model is low across all
policies. The meaningful comparison is the \emph{relative} improvement
between retention strategies at each budget level, which isolates the
effect of the retention policy from the LLM's intrinsic reasoning
ability.

%--------------------------------------------------------------------
\subsection{Jaccard Aggregator Comparison}

\begin{table}[h]
\centering
\caption{Median vs.\ Jaccard aggregator on \texttt{HeavyDuplication\_50}.}
\label{tab:jaccard}
\begin{tabular}{@{}rcccc@{}}
\toprule
$B/N$ & \multicolumn{2}{c}{Recall@10} & \multicolumn{2}{c}{Throughput (items/s)} \\
\cmidrule(lr){2-3}\cmidrule(lr){4-5}
      & Median & Jaccard & Median & Jaccard \\
\midrule
 1\% & 0.0193 & \textbf{0.0197} &  84,465 & 4,915 \\
 5\% & \textbf{0.0989} & 0.0948 &  77,602 & 1,279 \\
10\% & \textbf{0.1972} & 0.1918 &  68,415 &   523 \\
25\% & \textbf{0.4992} & 0.4660 &  48,118 &   196 \\
50\% & \textbf{1.0000} & \textbf{1.0000} &  41,190 &   126 \\
\bottomrule
\end{tabular}
\end{table}

The Median aggregator consistently outperforms Jaccard at medium budgets
(5\%--25\%). Jaccard shows a marginal advantage only at the 1\% budget
(0.0197 vs.\ 0.0193). More critically, Jaccard is \textbf{dramatically
slower}: at the 25\% budget it processes only 196~items/s compared to
48,118~items/s for Median---a $245\times$ slowdown. The inverted-index
lookups and per-candidate comparisons scale poorly as $|M|$ grows.

The Median aggregator is the recommended default. Jaccard's marginal
precision gains at extreme compression do not justify its computational
cost.

%--------------------------------------------------------------------
\subsection{JL Projection Memory--Recall Tradeoff}

In contrast to Sections~V-B and V-C, which define the budget as an
\emph{item count} $B$, this section holds the \emph{byte budget} fixed:
every row in a column shares the same total storage
$B_{\text{bytes}} = B \times d \times 4$~bytes (32-bit floats). A
$d{=}64$ system therefore stores $6\times$ more items than $d{=}384$
for the same memory cost.

\begin{table}[h]
\centering
\caption{JL projection tradeoff: recall under fixed memory footprint.}
\label{tab:jl}
\begin{tabular}{@{}rlrrr@{}}
\toprule
Budget & Policy & $d$ & Items & R@10 \\
\midrule
\multirow{3}{*}{1\%}
 & FIFO (baseline)     & 384 &   20,000 & 0.0181 \\
 & BucketOcc.          & 384 &   20,000 & 0.0193 \\
 & BucketOcc.\ + JL    &  64 &  120,000 & \textbf{0.0971} \\
\midrule
\multirow{3}{*}{5\%}
 & BucketOcc.          & 384 &  100,000 & 0.0989 \\
 & BucketOcc.\ + JL    & 128 &  300,000 & 0.2588 \\
 & BucketOcc.\ + JL    &  64 &  600,000 & \textbf{0.3893} \\
\midrule
\multirow{3}{*}{10\%}
 & BucketOcc.          & 384 &  200,000 & 0.1972 \\
 & BucketOcc.\ + JL    & 128 &  600,000 & 0.4679 \\
 & BucketOcc.\ + JL    &  64 & 1,200,000 & \textbf{0.5561} \\
\midrule
\multirow{3}{*}{25\%}
 & BucketOcc.          & 384 &  500,000 & 0.4992 \\
 & BucketOcc.\ + JL    & 128 & 1,500,000 & \textbf{0.6772} \\
 & BucketOcc.\ + JL    &  64 & 3,000,000$^*$ & 0.5561 \\
\midrule
\multirow{3}{*}{50\%}
 & BucketOcc.          & 384 & 1,000,000 & \textbf{1.0000} \\
 & BucketOcc.\ + JL    & 128 & 3,000,000$^*$ & 0.6772 \\
 & BucketOcc.\ + JL    &  64 & 6,000,000$^*$ & 0.5561 \\
\bottomrule
\multicolumn{5}{@{}l}{\footnotesize $^*$Capacity exceeds unique items in stream (1M).}
\end{tabular}
\end{table}

\textbf{Tight budgets ($<$15\%): JL dominates.}
At the 5\% float budget, $d{=}64$ achieves $3.94\times$ higher recall
than uncompressed BucketOccupancy (0.3893 vs.\ 0.0989). Storing
$6\times$ more lossy items drastically outweighs the precision lost
from compression.

\textbf{The ``distortion ceiling.''}
JL-projected policies hit hard recall caps regardless of available
memory: $d{=}64$ plateaus at 0.5561 and $d{=}128$ at 0.6772. Even when
capacity exceeds the number of unique items, geometric distortion
prevents perfect retrieval.

\textbf{Loose budgets ($>$25\%): no JL.}
Uncompressed $d{=}384$ vectors achieve perfect 1.0000 recall at 50\%
budget, where JL variants are permanently capped by their distortion
ceilings.

\textbf{Crossover at 25\%.}
At this budget, $d{=}128$ (0.6772) outperforms $d{=}64$ (0.5561)
because the $d{=}64$ ceiling has already been reached. Both JL
variants surpass the uncompressed baseline (0.4992).

\textbf{Practical recommendation:} Use JL projection ($d{=}64$) when
the memory budget is $<$15\% of the stream; use $d{=}128$ at 15--25\%;
avoid JL entirely when budget exceeds 25\%.

%------------------------------------------------------------------------
\section{Conclusion}

\subsection{Summary of Contributions}
\begin{enumerate}
\item \textbf{A novel $O(L)$ retention policy} using LSH bucket
      occupancy as an intrinsic novelty signal, avoiding the $O(|M|)$
      cost of exact deduplication while achieving comparable or superior
      recall. The auxiliary data structures add $<$1\% memory overhead
      beyond the embedding storage itself.
\item \textbf{Perfect recall under duplication:} On a 50\%-duplicate
      stream, BucketOccupancy achieves 1.0000 recall at 50\% memory
      budget where all baselines plateau at $\sim$0.69.
\item \textbf{Downstream reasoning improvement:} At 25\% memory budget,
      BucketOccupancy improves LLM-QA accuracy by 52.3\% over FIFO in
      the LoCoMo benchmark.
\item \textbf{Discovery of the JL ``distortion ceiling'':} Establishing
      clear operating regimes where dimensionality reduction helps
      ($<$15\% budget) and where it harms ($>$25\%).
\item \textbf{An honest negative result:} The policy underperforms FIFO
      under pure topic drift, establishing clear boundaries for
      applicability.
\end{enumerate}

\subsection{Deployment Guidance}

\begin{table}[h]
\centering
\caption{Deployment recommendations by scenario.}
\label{tab:deploy}
\begin{tabular}{@{}lll@{}}
\toprule
Scenario & Recommendation & Reason \\
\midrule
Redundancy-heavy streams & BucketOcc.\ ($d{=}384$) & Perfect recall at 50\% \\
Memory $<$ 15\% & BucketOcc.\ + JL ($d{=}64$) & $4\times$ recall via quantity \\
Memory 15--25\% & BucketOcc.\ + JL ($d{=}128$) & Best before ceiling \\
Memory $>$ 25\% & BucketOcc.\ ($d{=}384$) & Avoids distortion \\
Pure topic drift & FIFO & Recency $\approx$ relevance \\
Ultra-low latency & Reservoir & $\sim$12$\times$ faster \\
\bottomrule
\end{tabular}
\end{table}

\subsection{Limitations}
\begin{itemize}
\item Throughput of $\sim$75K items/s may be insufficient for extremely
      high-frequency streams ($>$1M items/s).
\item The Jaccard variant's $245\times$ throughput penalty limits
      practical utility despite marginal precision gains.
\item LoCoMo evaluation used a 7B parameter model; larger models may
      narrow the gap between policies.
\item The policy assumes a fixed hash function family; adaptive hashing
      under drift is unexplored.
\end{itemize}

\subsection{Future Work}
\begin{itemize}
\item \textbf{Adaptive $K$ and $L$:} Dynamically adjust hash
      granularity as the stream evolves.
\item \textbf{Hybrid drift detection:} Combine BucketOccupancy with a
      drift detector to switch to FIFO when drift is detected.
\item \textbf{Multi-modal streams:} Extend to image/audio embeddings
      beyond text.
\item \textbf{Production integration:} Embed as a memory-management
      layer in LangChain/LlamaIndex agent frameworks.
\end{itemize}

\section*{Acknowledgment}
The authors thank the NYU Courant Institute for computational resources
and the course staff of Algorithms for their guidance on the theoretical
framing.

\begin{thebibliography}{00}
\bibitem{b8} S. Muthukrishnan, ``Data streams: Algorithms and applications,'' \textit{Foundations and Trends in Theoretical Computer Science}, vol.~1, no.~2, pp. 117--236, 2005.
\bibitem{b9} M. Datar, A. Gionis, P. Indyk, and R. Motwani, ``Maintaining stream statistics over sliding windows,'' \textit{SIAM J. Comput.}, vol.~31, no.~6, pp. 1794--1813, 2002.
\bibitem{b10} J.~S. Vitter, ``Random sampling with a reservoir,'' \textit{ACM Trans. Math. Softw.}, vol.~11, no.~1, pp. 37--57, 1985.
\bibitem{b11} A. Abbas, K. Tirumala, D. Simig, S. Ganguli, and A.~S. Morcos, ``SemDeDup: Data-efficient learning at web-scale through semantic deduplication,'' in \textit{Proc. NeurIPS Workshop on Efficient Natural Language and Speech Processing}, 2023.
\bibitem{b12} D. Kraus, I. Carmel, and I. Keidar, ``Stream-LSH: Applying locality-sensitive hashing to stream data management,'' in \textit{Proc. IEEE BigData}, pp. 2444--2453, 2017.
\bibitem{b13} P. Indyk and R. Motwani, ``Approximate nearest neighbors: Towards removing the curse of dimensionality,'' in \textit{Proc. 30th ACM Symp. Theory of Computing (STOC)}, pp. 604--613, 1998.
\bibitem{b14} M.~S. Charikar, ``Similarity estimation techniques from rounding algorithms,'' in \textit{Proc. 34th ACM Symp. Theory of Computing (STOC)}, pp. 380--388, 2002.
\bibitem{b15} A. Blum, J. Hopcroft, and R. Kannan, \textit{Foundations of Data Science}. Cambridge University Press, 2020, ch.~2.
\bibitem{b16} W.~B. Johnson and J. Lindenstrauss, ``Extensions of Lipschitz mappings into a Hilbert space,'' \textit{Contemp. Math.}, vol.~26, pp. 189--206, 1984.
\bibitem{b17} P. Bajaj, D. Campos, N. Craswell, L. Deng, J. Gao, X. Liu, R. Majumder, A. McNamara, B. Mitra, T. Nguyen, M. Rosenberg, X. Song, A. Stoica, S. Tiwary, and T. Wang, ``MS MARCO: A human generated machine reading comprehension dataset,'' \textit{arXiv preprint arXiv:1611.09268}, 2016.
\end{thebibliography}
\end{document}
